\documentclass{article} % For LaTeX2e
\usepackage{iclr2026_conference,times}

\input{math_commands.tex}

\usepackage{hyperref}
\usepackage{url}
\usepackage{graphicx}
\usepackage{booktabs}
\usepackage{amsmath,amssymb}
\usepackage{microtype}
\usepackage{enumitem}
\usepackage{array}
\usepackage{fancyvrb}

\hypersetup{
  colorlinks=true,
  linkcolor=blue,
  citecolor=blue,
  urlcolor=blue
}

\iclrfinalcopy

\title{Supplementary Material for\\Staged Metric-Gated GRPO Fine-Tuning Pipeline for Visual Numeric Reasoning}
\author{Kunal Ghosh\\
Purdue University\\
West Lafayette, Indiana, USA\\
\texttt{ghosh178@purdue.edu}\\
\texttt{kg.aero@gmail.com}}

\newcolumntype{L}[1]{>{\raggedright\arraybackslash}p{#1}}

\begin{document}

\maketitle
\fancyhead{}
\lhead{}
\chead{}
\rhead{}

\section{Purpose of the Supplement}

This supplement expands the main paper in five directions:
\begin{enumerate}[leftmargin=1.2em]
\item the exact reward and evaluation formulas used by the method,
\item the controller rules and checkpoint-selection logic,
\item the runtime profile used by the reported runs,
\item the final full-evaluation protocol, and
\item the continuation semantics used for same-phase resume versus cross-phase warm start.
\end{enumerate}
All statements below are grounded in the following implementation modules:
\begin{itemize}[leftmargin=1.2em]
\item \texttt{staged\_rl/config.py}
\item \texttt{staged\_rl/rewarding.py}
\item \texttt{staged\_rl/evaluation.py}
\item \texttt{staged\_rl/controller.py}
\item \texttt{staged\_rl/checkpointing.py}
\item \texttt{staged\_rl/trainer\_runtime.py}
\item \texttt{rl\_gspo\_qwen2\_5vlm\_eval.py}
\end{itemize}

\section{GRPO Versus the Contribution of This Paper}

Our implementation uses TRL's \texttt{GRPOTrainer} as its optimization backend. The methodological contribution analyzed in the main paper is the fine-tuning pipeline built around that backend rather than a new GRPO objective:
\begin{itemize}[leftmargin=1.2em]
\item staged curriculum construction over filtered MathVista subsets,
\item a structured response contract,
\item metric-gated reward reweighting after checkpoint evaluation, and
\item alias-based checkpoint selection for continuation.
\end{itemize}
GRPO provides the policy-optimization backbone: the trainer samples multiple completions for a prompt, evaluates them under a reward function, and updates the policy from relative within-group performance rather than from PPO-style clipped ratios. The contribution in this work is the training system built around that optimizer: filtered curricula, metric-gated reward adaptation, checkpoint-aware continuation, and auditable transfer semantics.

\section{Full Reward Specification}

Let $c$ denote a completion, $g$ the gold answer, $m \in \{\mathrm{num}, \mathrm{mc}\}$ the answer mode, and $L$ the phase-specific completion budget. Here $\mathrm{num}$ denotes the numeric free-form path and $\mathrm{mc}$ the multi-choice path. The total per-sample reward is
\begin{equation}
R(c,g,m)=
w_f R_{\mathrm{format}} +
w_p R_{\mathrm{parseable}} +
w_h R_{\mathrm{finished}} +
w_c R_{\mathrm{correctness}} +
w_b R_{\mathrm{brevity}} +
w_t R_{\mathrm{tolerance}}.
\end{equation}

\paragraph{Tag and parsing predicates.}
The implementation uses:
\begin{align}
\mathrm{ReasoningTag}(c) &= \mathbf{1}[\text{exactly one reasoning block exists}], \\
\mathrm{SolutionTag}(c) &= \mathbf{1}[\text{exactly one solution block exists}], \\
\mathrm{Malformed}(c,m) &=
\begin{cases}
\mathbf{1}[\text{numeric extraction or parsing fails}], & m=\mathrm{num},\\
\mathbf{1}[\text{option extraction fails}], & m=\mathrm{mc},
\end{cases}\\
\mathrm{Finished}(c,m) &=
\begin{cases}
1, & \mathrm{ReasoningTag}(c)=1 \wedge \mathrm{SolutionTag}(c)=1 \\
   & \wedge\ \mathrm{Malformed}(c,m)=0,\\
0, & \text{otherwise}.
\end{cases}
\end{align}

\paragraph{Reward components.}
The six reward components are:
\begin{align}
R_{\mathrm{format}}(c,m) &= \mathrm{ReasoningTag}(c) + \mathrm{SolutionTag}(c)
- 0.5\,\mathrm{Malformed}(c,m),\\
R_{\mathrm{parseable}}(c,m) &=
\begin{cases}
\mathbf{1}[\text{numeric parsing succeeds}], & m=\mathrm{num},\\
\mathbf{1}[\mathrm{Malformed}(c,m)=0], & m=\mathrm{mc},
\end{cases}\\
R_{\mathrm{finished}}(c,m) &=
\begin{cases}
1, & \mathrm{Finished}(c,m)=1 \wedge \mathrm{tok}(c)<L,\\
-1, & \mathrm{tok}(c)\ge L \wedge \mathrm{Finished}(c,m)=0,\\
0, & \text{otherwise},
\end{cases}\\
R_{\mathrm{correctness}}(c,g,m) &= \mathbf{1}[\text{exact match}],\\
R_{\mathrm{brevity}}(c) &=
\begin{cases}
-1, & \mathrm{tok}(c)\ge L,\\
1, & \mathrm{tok}(c)/L \le 0.55,\\
0.5, & 0.55 < \mathrm{tok}(c)/L \le 0.75,\\
0, & 0.75 < \mathrm{tok}(c)/L \le 0.90,\\
-0.5, & 0.90 < \mathrm{tok}(c)/L < 1,
\end{cases}\\
R_{\mathrm{tolerance}}(c,g,m) &=
\begin{cases}
1, & \text{tolerance is enabled, exact match is false,}\\
   & \text{and the tolerance check holds},\\
0, & \text{otherwise}.
\end{cases}
\end{align}

Tolerance reward is active in Phase C. Exact match suppresses tolerance reward, so the tolerance component acts strictly as a denser near-miss signal rather than an additional bonus for already-correct answers.

\paragraph{Phase-initialized weights.}
The configured initial weights for the reported training phases are manual curriculum defaults rather than values learned by a meta-optimizer:
\begin{center}
\small
\begin{tabular}{lcccccc}
\toprule
Phase & Format & Parseable & Finished & Correctness & Brevity & Tolerance \\
\midrule
A & 1.0 & 1.0 & 1.5 & 2.0 & 0.25 & 0.0 \\
B & 0.75 & 0.75 & 1.0 & 4.0 & 0.20 & 0.0 \\
C & 0.50 & 0.50 & 0.75 & 5.0 & 0.20 & 1.0 \\
\bottomrule
\end{tabular}
\end{center}

\section{Evaluation Metrics and Checkpoint Scores}

For each prompt $i$, the evaluator samples one or more completions $c_{i,j}$. The final reported evaluation uses one completion per prompt. The evaluator computes:
\begin{align}
\mathrm{EM} &= \frac{1}{N}\sum_{i=1}^N \mathbf{1}[\text{first sampled completion matches exactly}],\\
\mathrm{TolAcc} &= \frac{1}{N}\sum_{i=1}^N \mathbf{1}[\text{first sampled completion is correct under tolerance}],\\
\mathrm{ParseRate} &= \frac{1}{M}\sum_{i,j}\mathbf{1}[\text{sample } (i,j) \text{ is parseable}],\\
\mathrm{MalfRate} &= \frac{1}{M}\sum_{i,j}\mathbf{1}[\text{sample } (i,j) \text{ is malformed}],\\
\mathrm{TruncRate} &= \frac{1}{M}\sum_{i,j}\mathbf{1}[\text{sample } (i,j) \text{ is truncated}].
\end{align}
It also computes tag-compliance rates, completion-token averages, repetition rate, sampled-answer diversity, and mean logged reward components.

The checkpoint registry then defines three scores:
\begin{align}
S_{\mathrm{struct}} &= 0.35\,\mathrm{ParseRate} + 0.20\,\mathrm{SolutionTag}
+ 0.20\,\mathrm{ReasoningTag} \nonumber\\
&\quad - 0.15\,\mathrm{MalfRate} - 0.10\,\mathrm{TruncRate},\\
S_{\mathrm{corr}} &= 0.70\,\mathrm{EM} + 0.20\,\mathrm{TolAcc} + 0.10\,\mathrm{ParseRate},\\
S_{\mathrm{comp}} &= 0.45\,\mathrm{EM} + 0.15\,\mathrm{TolAcc} + 0.15\,\mathrm{ParseRate}
+ 0.10\,\mathrm{SolutionTag} \nonumber\\
&\quad + 0.05\,\mathrm{ReasoningTag}
- 0.05\,\mathrm{MalfRate} - 0.05\,\mathrm{TruncRate}.
\end{align}
These scores back the aliases \texttt{best\_structure}, \texttt{best\_correctness}, and \texttt{best\_composite}. The \texttt{latest} alias is purely chronological.

\section{Controller Rules and Thresholds}

\begin{table}[t]
\centering
\small
\caption{Controller rules implemented in \texttt{staged\_rl/controller.py}.}
\label{tab:controller-rules}
\begin{tabular}{L{0.19\linewidth} L{0.26\linewidth} L{0.41\linewidth}}
\toprule
Rule & Trigger & Action \\
\midrule
Parseable guard & parseable answer rate $< 0.85$ & Raise parseable reward to at least $0.75$ \\
Format guard & solution tag $< 0.90$ or reasoning tag $< 0.88$ or malformed rate $> 0.10$ & Raise format reward to at least $0.75$ \\
Finish guard & truncation rate $> 0.15$ or average completion tokens exceed $0.85L$ & Add $0.25$ to finished reward \\
Stable structure test & parseable $\ge 0.92$, solution $\ge 0.95$, reasoning $\ge 0.93$, malformed $\le 0.05$, truncation $\le 0.08$ & Enable plateau-based correctness escalation once enough history exists \\
Correctness escalation & Stable structure holds over a two-checkpoint window and exact-match gain is below $0.02$ & Add $0.50$ to correctness reward \\
\bottomrule
\end{tabular}
\end{table}

The controller rules are intentionally concise and inspectable. The main paper does not rely on a fixed count of controller firings because the final primary comparison is based on selected checkpoints reevaluated on the full numeric evaluation subset.

\section{Runtime Profile and Final Evaluation Protocol}

The reported training and final evaluation share the same memory-constrained adapter and generation budget:
\begin{itemize}[leftmargin=1.2em]
\item max sequence length $1280$,
\item image size $336$,
\item LoRA rank $8$ with LoRA alpha $8$,
\item max prompt length $320$,
\item max completion length $64$,
\item two sampled generations per prompt during GRPO fine-tuning,
\item one sampled completion per prompt during final evaluation, and
\item full final evaluation with no per-subset example cap.
\end{itemize}

\paragraph{Runtime configuration.}
The same runtime profile can be read directly in the style of the code paths that consume these settings:

\begin{Verbatim}[fontsize=\scriptsize,frame=single,framesep=3mm]
# Model load / LoRA adapter profile
FastVisionModel.from_pretrained(
    model_name="unsloth/Qwen2.5-VL-7B-Instruct",
    max_seq_length=1280,
    load_in_4bit=True,
    fast_inference=True,
    gpu_memory_utilization=0.65,
    max_lora_rank=8,
    compilation_config={"level": 3, "cudagraph_mode": "PIECEWISE"},
)

FastVisionModel.get_peft_model(
    model,
    r=8,
    lora_alpha=8,
    finetune_vision_layers=False,
    finetune_language_layers=True,
    finetune_attention_modules=True,
    finetune_mlp_modules=True,
)
\end{Verbatim}

\begin{Verbatim}[fontsize=\scriptsize,frame=single,framesep=3mm]
# Trainer / final evaluation profile
GRPOConfig(
    gradient_accumulation_steps=4,
    num_generations=2,
    max_prompt_length=320,
    max_completion_length=64,
    importance_sampling_level="sequence",
    loss_type="dr_grpo",
)

SamplingParams(
    temperature=1.0,
    top_k=50,
    max_tokens=64,
    n=1,
)

EvalConfig(
    num_samples_per_prompt=1,
    max_eval_examples_per_subset=None,
)
\end{Verbatim}

\paragraph{Final evaluation target set.}
The final comparison evaluates four targets:
\begin{center}
\small
\begin{tabular}{ll}
\toprule
Label & Checkpoint \\
\midrule
Baseline & \texttt{baseline\_rank8} \\
Phase A & \texttt{phase\_a/checkpoint-237} \\
Phase B & \texttt{phase\_b/checkpoint-180} \\
Phase C & \texttt{phase\_c/checkpoint-240} \\
\bottomrule
\end{tabular}
\end{center}
The primary subset is \texttt{eval\_overall\_numeric} from MathVista \texttt{testmini}. Its size is $456$ examples. The diagnostic stage subsets contain $195$ Stage 1 examples, $287$ Stage 2 examples, and $119$ Stage 3 examples.

\section{Full Evaluation Results}

\begin{table}[t]
\centering
\small
\caption{Final primary evaluation on the full \texttt{eval\_overall\_numeric} subset.}
\label{tab:supp-headline}
\begin{tabular}{llrrrrrrr}
\toprule
Phase & Checkpoint & Exact & Tol & Parse & AvgR & Struct & Corr & Comp \\
\midrule
Baseline & \texttt{baseline\_rank8} & 0.1798 & 0.1820 & 0.2544 & -0.0880 & 0.0285 & 0.1877 & 0.1163 \\
Phase A & \texttt{checkpoint-237} & 0.4803 & 0.4868 & 0.9890 & 6.3857 & 0.7421 & 0.5325 & 0.5859 \\
Phase B & \texttt{checkpoint-180} & 0.4978 & 0.5066 & 0.9781 & 5.8754 & 0.7371 & 0.5476 & 0.5946 \\
Phase C & \texttt{checkpoint-240} & 0.4737 & 0.4846 & 0.9890 & 6.0361 & 0.7436 & 0.5274 & 0.5833 \\
\bottomrule
\end{tabular}
\end{table}

\begin{table}[t]
\centering
\scriptsize
\caption{Full stage-wise diagnostic results on MathVista \texttt{testmini}.}
\label{tab:supp-stage}
\begin{tabular}{llrrrrrr}
\toprule
Phase & Eval subset & Exact & Tol & Parse & Malf & Trunc & AvgR \\
\midrule
Baseline & Stage 1 & 0.1641 & 0.1641 & 0.2872 & 0.7128 & 0.6923 & 0.1981 \\
Baseline & Stage 2 & 0.1220 & 0.1220 & 0.1672 & 0.8328 & 0.8293 & -0.7321 \\
Baseline & Stage 3 & 0.0924 & 0.0924 & 0.1681 & 0.8319 & 0.8151 & -0.8057 \\
Phase A & Stage 1 & 0.4615 & 0.4615 & 0.9897 & 0.0103 & 0.0051 & 6.3417 \\
Phase A & Stage 2 & 0.4321 & 0.4460 & 0.9895 & 0.0105 & 0.0000 & 6.2666 \\
Phase A & Stage 3 & 0.3950 & 0.4034 & 0.9832 & 0.0168 & 0.0084 & 6.1261 \\
Phase B & Stage 1 & 0.4718 & 0.4718 & 0.9846 & 0.0154 & 0.0000 & 5.7786 \\
Phase B & Stage 2 & 0.4774 & 0.4983 & 0.9721 & 0.0279 & 0.0035 & 5.7652 \\
Phase B & Stage 3 & 0.3866 & 0.4118 & 0.9664 & 0.0336 & 0.0000 & 5.3517 \\
Phase C & Stage 1 & 0.4564 & 0.4564 & 0.9692 & 0.0308 & 0.0000 & 5.8790 \\
Phase C & Stage 2 & 0.4321 & 0.4530 & 0.9861 & 0.0139 & 0.0000 & 5.8103 \\
Phase C & Stage 3 & 0.4202 & 0.4286 & 0.9832 & 0.0168 & 0.0000 & 5.7294 \\
\bottomrule
\end{tabular}
\end{table}

\section{Transfer and Continuation Semantics}

The code distinguishes three mechanisms that can be conflated in informal descriptions:
\begin{enumerate}[leftmargin=1.2em]
\item \textbf{Artifact transfer.} A checkpoint directory is made available to a later run as an input artifact.
\item \textbf{Warm start.} A fresh base model is loaded, a fresh LoRA adapter is attached, and adapter weights are copied from the selected checkpoint into that new adapter.
\item \textbf{Same-phase resume.} Full trainer state is restored from a checkpoint, including optimizer state, scheduler state, trainer progress, and controller state.
\end{enumerate}

This distinction is important for interpreting the reported checkpoints. Cross-phase continuation uses selected adapter weights as a warm start. Same-phase restarts can use exact trainer resume if the selected checkpoint belongs to the current phase and the full trainer state is available.

\section{Additional Scope Notes}

The implementation contains preliminary support for a multi-choice branch and an optional robustness stage, but neither path provides the main evidence used in the paper. Similarly, the runner logs when multiple GPUs are visible but does not automatically switch to DDP or tensor parallelism. The primary claims are therefore:
\begin{itemize}[leftmargin=1.2em]
\item about staged GRPO fine-tuning pipeline dynamics for numeric visual reasoning,
\item about metric-gated control under non-monotonic checkpoint trajectories, and
\item about making a memory-constrained GRPO fine-tuning pipeline more auditable and reproducible.
\end{itemize}

\end{document}
